% Source: results/v6/char_vs_word_bm25_v6.csv, results/v6/conformal_splits_summary_v6.csv, results/v6/multihop_independent_v6.csv, results/v6/coco_native_visual_v3.csv
\begin{abstract}
Multimodal Retrieval-Augmented Generation (MM-RAG) frequently indexes precomputed textual proxies (OCR transcripts, machine captions) in place of raw pixels. Across a standardized 1,626-item benchmark under 50/50 Dev/Test splits ($N=306$ held-out probes, 813 test items), we investigate when text proxies suffice and when native-visual representations are indispensable. Contrary to common assumptions, character 3--5gram BM25 strictly dominates word BM25 across all verbatim optical OCR regimes (Clean: 0.9436 vs.\ 0.8371; Degraded: 0.8248 vs.\ 0.6972; Tesseract 5.5: 0.8410 vs.\ 0.7361; all $p_{\text{Holm}} < 10^{-4}$), reversing only under entity-substituting light paraphrasing (0.2427 vs.\ 0.4334). On ground-truth text, dense bi-encoder failure on verbatim spans (MRR 0.2100) is governed by query formulation mismatch; answerable questions restore BGE-Base MRR to 0.5197 (Hit@10 78.8\%). Across 200 cluster-stratified splits, conformal cascade calibration under marginal loss rigorously bounds test miss risk below 0.05 (mean test risk 0.0491 $\le 0.05$) at 78.4\% escalation savings. On MS-COCO ($N=1{,}000$), native visual retrieval (SigLIP Base Patch16 MRR = 0.6071) decisively outperforms caption proxies (0.3728, $p < 10^{-4}$). In multi-hop reasoning, independent pseudo-relevance feedback matches generative LLM bridging (Hit@10 0.640 vs.\ 0.640, exact McNemar $p = 1.000$) with zero API cost.
\end{abstract}
